%%%%%%%%%%%%%%%%%%%%%%%%%%%%%%%%%%%%%%%%%
% Short Sectioned Assignment LaTeX Template Version 1.0 (5/5/12)
% This template has been downloaded from: http://www.LaTeXTemplates.com
% Original author:  Frits Wenneker (http://www.howtotex.com)
% License: CC BY-NC-SA 3.0 (http://creativecommons.org/licenses/by-nc-sa/3.0/)
%%%%%%%%%%%%%%%%%%%%%%%%%%%%%%%%%%%%%%%%%

% \documentclass[paper=a4, fontsize=11pt]{scrartcl} % A4 paper and 11pt font size
\documentclass[11pt, a4paper]{book}
\usepackage[T1]{fontenc} % Use 8-bit encoding that has 256 glyphs
\usepackage[utf8]{inputenc}
\usepackage{fourier} % Use the Adobe Utopia font for the document - comment this line to return to the LaTeX default
\usepackage{listings} % para insertar código con formato similar al editor
\usepackage[spanish, es-tabla]{babel} % Selecciona el español para palabras introducidas automáticamente, p.ej. "septiembre" en la fecha y especifica que se use la palabra Tabla en vez de Cuadro
\usepackage{url} % ,href} %para incluir URLs e hipervínculos dentro del texto (aunque hay que instalar href)
\usepackage{graphics,graphicx, float} %para incluir imágenes y colocarlas
\usepackage[gen]{eurosym} %para incluir el símbolo del euro
\usepackage{cite} %para incluir citas del archivo <nombre>.bib
\usepackage{enumerate}
\usepackage{hyperref}
\usepackage{graphicx}
\usepackage{tabularx}
\usepackage{booktabs}

\usepackage[table,xcdraw]{xcolor}
\hypersetup{
	colorlinks=true,	% false: boxed links; true: colored links
	linkcolor=black,	% color of internal links
	urlcolor=cyan		% color of external links
}
\renewcommand{\familydefault}{\sfdefault}
\usepackage{fancyhdr} % Custom headers and footers
\pagestyle{fancyplain} % Makes all pages in the document conform to the custom headers and footers
\fancyhead[L]{} % Empty left header
\fancyhead[C]{} % Empty center header
\fancyhead[R]{Me Llamo Así} % My name
\fancyfoot[L]{} % Empty left footer
\fancyfoot[C]{} % Empty center footer
\fancyfoot[R]{\thepage} % Page numbering for right footer
%\renewcommand{\headrulewidth}{0pt} % Remove header underlines
\renewcommand{\footrulewidth}{0pt} % Remove footer underlines
\setlength{\headheight}{13.6pt} % Customize the height of the header

\usepackage{titlesec, blindtext, color}
\definecolor{gray75}{gray}{0.75}
\newcommand{\hsp}{\hspace{20pt}}
\titleformat{\chapter}[hang]{\Huge\bfseries}{\thechapter\hsp\textcolor{gray75}{|}\hsp}{0pt}{\Huge\bfseries}
\setcounter{secnumdepth}{4}
\usepackage[Lenny]{fncychap}


\begin{document}

	% Plantilla portada UGR
	\begin{titlepage}
\newlength{\centeroffset}
\setlength{\centeroffset}{-0.5\oddsidemargin}
\addtolength{\centeroffset}{0.5\evensidemargin}
\thispagestyle{empty}

\noindent\hspace*{\centeroffset}\begin{minipage}{\textwidth}

\centering
\includegraphics[width=0.9\textwidth]{logos/logo_ugr.jpg}\\[1.4cm]

\textsc{ \Large TRABAJO FIN DE GRADO\\[0.2cm]}
\textsc{ GRADO EN INGENIERIA INFORMATICA}\\[1cm]

{\Huge\bfseries Título \\}
\noindent\rule[-1ex]{\textwidth}{3pt}\\[3.5ex]
{\large\bfseries Subtítulo }
\end{minipage}

\vspace{2.5cm}
\noindent\hspace*{\centeroffset}
\begin{minipage}{\textwidth}
\centering

\textbf{Autor}\\ {Estudiante}\\[2.5ex]
\textbf{Director}\\ {Tutor(a)(es)}\\[2cm]
\includegraphics[width=0.3\textwidth]{logos/etsiit_logo.png}\\[0.1cm]
\textsc{Escuela Técnica Superior de Ingenierías Informática y de Telecomunicación}\\
\textsc{---}\\
Granada, Junio de 201x
\end{minipage}
\end{titlepage}


	% Plantilla prefacio UGR
	\thispagestyle{empty}

\begin{center}
{\large\bfseries Título \\ Subtítulo }\\
\end{center}
\begin{center}
Nombre Del Estudiante\\
\end{center}

%\vspace{0.7cm}

\vspace{0.5cm}
\noindent\textbf{Palabras clave}: \textit{software libre}
\vspace{0.7cm}

\noindent\textbf{Resumen}\\
	

\cleardoublepage

\begin{center}
	{\large\bfseries Same, but in English}\\
\end{center}
\begin{center}
	Student's name\\
\end{center}
\vspace{0.5cm}
\noindent\textbf{Keywords}: \textit{open source}, \textit{floss}
\vspace{0.7cm}

\noindent\textbf{Abstract}\\


\cleardoublepage

\thispagestyle{empty}

\noindent\rule[-1ex]{\textwidth}{2pt}\\[4.5ex]

D. \textbf{Tutora/e(s)}, Profesor(a) del ...

\vspace{0.5cm}

\textbf{Informo:}

\vspace{0.5cm}

Que el presente trabajo, titulado \textit{\textbf{Chief}},
ha sido realizado bajo mi supervisión por \textbf{Estudiante}, y autorizo la defensa de dicho trabajo ante el tribunal
que corresponda.

\vspace{0.5cm}

Y para que conste, expiden y firman el presente informe en Granada a Junio de 2018.

\vspace{1cm}

\textbf{El/la director(a)/es: }

\vspace{5cm}

\noindent \textbf{(nombre completo tutor/a/es)}

\chapter*{Agradecimientos}






	% Índice de contenidos
	\newpage
	\tableofcontents

	% Índice de imágenes y tablas
	\newpage
	\listoffigures

	% Si hay suficientes se incluirá dicho índice
	\listoftables 
	\newpage

	% Introducción 
	\chapter{Introducción}

Este proyecto es software libre, y está liberado con la licencia \cite{gplv3}.

	% Descripción del problema y hasta donde se llega
	\chapter{Descripción del problema}



	% Estado del arte
	% 	1. Crítica al estado del arte
	% 	2. Propuesta
	\chapter{Estado del arte}

El software libre y sus licencias \cite{gplv3} ha permitido llevar a cabo una expansión del
aprendizaje de la informática sin precedentes.

	
	\chapter{Planificación}

\section{Metodología utilizada}


\section{Temporización}

\section{Seguimiento del desarrollo}


	% Análisis del problema
	% 1. Análisis de requisitos
	% 2. Análisis de las soluciones
	% 3. Solucion propuesta
	% 4. Análisis de seguridad
	\chapter{Análisis del problema}

\section{Análisis del prototipo}
 
Como se ha mencionado previamente en el apartado de descripción del problema, el objetivo principal del proyecto es dotar a GenoMus de una implementación de carácter profesional a prueba de futuro que sea mantenible por perfiles no técnicos.

El prototipo de GenoMus está implementado en javascript puro para ejecución sobre NodeJS. La sintaxis utilizada busca la utilización del "dialecto" funcional de javascript. Tanto el modelo de representación de datos como el modelo de cómputo siguen generalmente patrones de diseño de programación funcional. 

\subsection{Ventajas del planteamiento funcional del código en JS.}
Es bonito

\subsection{Placebos del JS funcional.} \ \
Javascript nos proporciona azúcar sintáctico sobre la creación de objetos dinámicos que nos da la posiblidad de implementar entidades como clausuras o funciones de orden superior. Sin embargo estas construcciones que tan bien nos entran por los ojos no están acompañadas de un respaldo acorde en el entorno de ejecución de js. Ninguno de los principales motores de js incluye algún tipo de análisis de los objetos funcionales para su optimización en grupo, como procedimientos de composición de funciones o un motor de reducción de grafos de dependencias computacionales.

La composición de funciones es la característica cuya ausencia puede haber condicionado el desarrollo del prototipo. El prototipo implementa un procedimiento de almacenamiento de árboles funcionales a través de cadenas de texto. Las diferentes construcciones funcionales se almacenan como código javascript, de cuya ejecución se obtiene la evaluación del árbol. Este tipo de metaprogramación aumenta el cómputo necesario para evaluar árboles de funciones, ya que se introduce una etapa de construcción de código y una etapa de parseo de este validación y posterior ejecución.

\section{Propuestas de implementación}

Tras analizar el prototipo se propone la implementación de las funcionalidades y críticas de GenoMus en C++. Este subconjunto de las funcionalidades se desarrollará bajo el nombre de genomus-core. Se propone reutilizar las funcionalidades secundarias del prototipo de javascript. La totalidad del paquete se desarrollará bajo el nombre de genomus-js.

Así, la base de código se dividirá en tres áreas:

\begin{enumerate} 
    \item Una librería de C++ que implementa el modelo de datos y de cómputo de GenoMus: genomus-core\cite{genomus-core}.
    \item Un módulo nativo de NodeJs (cita) que exponga una API de interación entre el runtime de JS y la biblioteca de C++: genomus-core-js\cite{genomus-core-js}.
    \item Un módulo de npm que encapsule lo previamente enunciado y reutilize la lógica de I/O del prototipo: genomus-js.
\end{enumerate}

	% Desarrollo bajo sprints: 
	% 	1. Permitir registros y login de usuarios
	% 	2. Desarrollo del sistema de incidencias
	% 	3. Desarrollo del sistema de denuncias administrativas y accidentes
	% 	4. Desarrollo del sistema de croquis
	%   5. Instalación de la aplicación de manera automática
	\chapter{Implementación}

Tras analizar el prototipo se propone la implementación de las funcionalidades y críticas de GenoMus en C++. Este subconjunto de las funcionalidades se desarrollará bajo el nombre de genomus-core. Se propone reutilizar las funcionalidades secundarias del prototipo de javascript. La totalidad del paquete se desarrollará bajo el nombre de genomus-js.

Así, la base de código se dividirá en tres áreas: (???)

\begin{enumerate} 
    \item Una librería de C++ que implementa el modelo de datos y de cómputo de GenoMus: genomus-core\cite{genomus-core}.
    \item Un módulo nativo de NodeJs (cita) que exponga una API de interación entre el runtime de JS y la biblioteca de C++: genomus-core-js\cite{genomus-core-js}.
    \item Un módulo de npm que encapsule lo previamente enunciado y reutilize la lógica de I/O del prototipo: genomus-js.
\end{enumerate}

\section{genomus-core}

\textbf{genomus-core} implementa el modelo de datos y de cómputo de GenoMus. 

\subsection{Modelo de datos}

Fácil

\subsection{Modelo de cómputo}
El modelo de cómputo de GenoMus está intrínsecamente basado en el paradigma de programación funcional. Esto hace que una implementación en C++ no sea trivial por diferentes motivos:

\begin{itemize}
    \item \textbf{Funciones dinámicas en tiempo de ejecución.} Las funciones de cómputo disponibles deben ser declarables en tiempo de ejecución. Ya que las funciones disponibles al modelo de cómputo son dependientes del estado del programa, es ideal que estas se puedan declarar dinámicamente en tiempo de ejecución. En la actual implementación, se proporciona una serie de funciones instanciadas en tiempo de carga.
    
    \item \textbf{Polimorfismo funcional.} Necesitamos construcciones que doten al modelo de cómputo de polimorfismo declarado en tiempo de ejecución. Debe ser posible no solo declarar nuevas funciones en tiempo de ejecución, sino también declarar nuevos tipos a usar como parámetros o salidas. Esto nos quita de la ecuación casi todos los constructos que C++ nos ofrece para conseguir polimorfismo dinámico, ya que el dominio del polimorfismo debe ser conocido en tiempo de compilación en C++. (cita requerida)
    
    \item \textbf{Árboles funcionales declarables y evaluables.} Toda instancia de cómputo debe no solo ser computable sino también almacenable. La instancia de árbol funcional debe existir independientemente de su evaluación. Así, una función de cómputo será un objeto invocable o functor cuya invocación contribuirá a la declaración del árbol funcional en memoria.
    
    (... ejemplo de código, diagrama de memoria(?))

    Esto, además de ser necesario para el funcionamiento del software, posibilita el diseño de algoritmos de cómputo sobre el árbol funcional más allá del típico recorrido por profundidad. (cita requerida - por ejemplo haskell graph reduction machine)
    
    \item \textbf{Funciones enumerables en tiempo de ejecución.} Las funciones de cómputo disponibles en tiempo de ejecución deben ser enumerables y referenciables según los diferentes métodos de acceso(cita a donde se explican las cosas codificadas/decodificadas).
\end{itemize}

Todas estas características son bien incompatibles con la metodología de desarrollo orientado a objetos estándar de C++ o bien incompatibles con la máxima del proyecto de ser asequible para programadores amateur, si pretendemos seguir las buenas prácticas de C++. Por este motivo se toma la decisión de alejarse de la práctica principal de polimorfismo funcional y de datos definida en los últimos estándares de C++ (C++11, C++17(cita requerida)) para realizar una implementación mixta entre orientada a objetos clásica y orientada a prototipos(cita requerida). 

Lo que nos aporta conceptualmente el diseño orientado a prototipos es la posibilidad de manejar parcialmente el tipado de nuestros objetos en tiempo de ejecución. Por otra parte, la principal desventaja que nos trae es la existencia de errores de tipado no conocidos por el compilador, lo cual es natural cuando lo que buscamos es precisamente declarar funciones de tipado arbitrario en tiempo de ejecución.

Así, se nos presenta un compromiso entre la seguridad del código y la legibilidad de este. La implementación propuesta busca maximizar la legibilidad y flexibilidad del código minimizando la cantidad de comprobaciones de tipado en tiempo de ejecución. Además, se plantea mitigar el riesgo del tipado dinámico mediante la creación de tests de integración.

\section{genomus-core-js}


	% Presupuesto

	% Conclusiones
	\chapter{Conclusiones y trabajos futuros}



	% Trabajos futuros


	
	\newpage
	\bibliography{bibliografia}
	\bibliographystyle{plain}
	
\end{document}

